% !TEX program = pdflatex
\documentclass[11pt]{article}
\setlength{\textwidth}{6.5in}
\setlength{\oddsidemargin}{0pt}
\setlength{\evensidemargin}{0pt}
\setlength{\textheight}{9in}
\setlength{\topmargin}{0pt}
\addtolength{\topmargin}{-\headheight}
\addtolength{\topmargin}{-\headsep}
\ifdefined\pdfpagewidth \pdfpagewidth=8.5in \pdfpageheight=11in \fi
\ifdefined\pagewidth \pagewidth=8.5in \pageheight=11in \fi
\usepackage{amsmath,amssymb}
\usepackage{array}
\usepackage{booktabs}
\usepackage[tableposition=top]{caption}
\usepackage{longtable}
\usepackage{enumitem}
\usepackage{tikz}
\usepackage{pgfplots}
\pgfplotsset{compat=1.18}
\emergencystretch=3em

% ---- notation ---------------------------------------------------------------
\newcommand{\um}{\,$\mu$m}
\newcommand{\kms}{km\,s$^{-1}$}
\newcommand{\flux}{erg\,cm$^{-2}$\,s$^{-1}$}
\newcommand{\mA}{m\AA}
\newcommand{\sci}[2]{#1\times10^{#2}}
\newcommand{\ion}[2]{#1\,\textsc{#2}}
\newcommand{\nuc}[2]{$^{#1}$#2}
\newcommand{\lxl}{$L_X/L_\mathrm{bol}$}

% ---- front matter -----------------------------------------------------------
\title{Lines in the Blank Sky: a Search of the SDSS Sky Fibres\\ for Laser Emission, and the Galaxies Behind Them}
\author{Charles C. Norton}
\date{}

\begin{document}
\maketitle

\begin{abstract}
Every BOSS and eBOSS plate of the Sloan Digital Sky Survey places some eighty fibres on
positions with no catalogued source, to measure the sky. A point source of line emission
too faint in broad bands for the imaging that chose those positions, a transmitter on
its own or beside a star too faint to catalogue, would appear in one of them as an
unresolved emission line that the plate's other sky fibres lack. We search the
$381{,}943$ sky-fibre spectra of the $3853$ plates of DR17, from $3600$ to
$10400$\,\AA\ at a resolution near $2000$, for such lines: matched-filter peaks at
$6\sigma$, off the airglow lines, as narrow as the instrument and absent from the plate's
other sky fibres. Absorption-like dips of the same shape, which no transmitter makes, are
put through every test the peaks are and serve as the null. A line that one fibre holds
at one wavelength at other positions on other plates is the detector's: $762$ of $3217$
clean peaks and $1813$ of $4108$ dips. Two or more peaks at one redshift, at the flux
ratios a galaxy's spectrum fixes, mark $156$ fibres, where the same test finds $4$ among
the dips and $12\pm4$ among shuffled peaks: $154$ galaxies at $z=0.02$ to $1.07$, six of
them seen twice on different plates, $122$ with a Legacy Surveys DR10 source within
$1.5''$ of median $r=23.1$, and two fibres of Galactic nebular lines. At $8\sigma$, with
the spectra that hold clean lines of both signs set aside, $485$ peaks and $336$ dips
remain. Legacy Surveys sources lie within $1.5''$ of $38$ per cent of the peaks against
$4.4$ per cent at positions $30''$ north, and of the dips less often than at those
positions, so about $160$ of the peaks are single lines of galaxies, and the other $323$
fall short of the $336$ dips. Read in each exposure, $359$ peaks and $241$ dips hold
through the frames, and net of their counterparts the $202$ persistent peaks fall short of
the $253$ persistent dips. The line that recurs most often at one sky position, at
$4732.6$\,\AA\ on $54$ plate-nights over seven years, is resolved and alone in the stack
of the $94$ spectra there: Ly$\alpha$ of a galaxy at $z=2.893$. The median line flux at
$8\sigma$, $8\times10^{-17}$\,\flux, is that of a transmitter aimed at Earth through a
$10$\,m aperture at $37$\,W from $10$\,pc and $3.7$\,kW from $100$\,pc, and with $0.68$
of such lines recovered, the $0.093$\,deg$^2$ the fibres cover bound such transmitters to
fewer than about $850$ per square degree.
\end{abstract}

% =============================================================================
\section{Introduction}\label{sec:intro}
% =============================================================================

Spectroscopic searches for laser emission look for an unresolved line in the spectrum
of a star~\cite{ReinesMarcy2002,TellisMarcy2017}, on the reasoning that a transmitter
lives near one. A transmitter need not: one in free space, or beside a star too faint
for any catalogue, has no spectrum of its own taken unless a fibre happens to fall on
it. Large fibre surveys do place fibres on empty sky, to measure the sky they subtract
from every other spectrum, and the Sloan Digital Sky Survey's BOSS~\cite{Dawson2013} and
eBOSS~\cite{Dawson2016} spectrographs~\cite{Smee2013} recorded some eighty such
spectra on each of thousands of plates. Their positions were chosen to hold nothing in
the survey's imaging, so a source there is fainter in broad bands than that imaging
reaches, and a line from it stands alone.

The same spectra catch what the imaging missed. A galaxy whose light is mostly in a few
emission lines is faint in broad bands and bright in its lines, and a sky fibre that
lands on one records them at one redshift; such lines are the commonest astrophysical
source a laser search of blank sky must recognise. This paper searches every sky fibre
of DR17's BOSS and eBOSS plates for unresolved emission lines, separates the galaxies by
their redshifts and by the sources the imaging holds at the fibres, and measures the rest,
in the combined spectra and in each exposure, against the same search for dips, put
through the same tests.

\subsection{Contributions}
\begin{enumerate}[leftmargin=*,itemsep=1pt]
\item A search of $381{,}943$ sky-fibre spectra on $3853$ plates for unresolved emission
lines absent from the plate's other sky fibres, with absorption-like dips put through
every test as the null.
\item A catalogue of the $154$ galaxies whose lines the search finds at one redshift,
with their Legacy Surveys counterparts.
\item A bound on transmitters in blank sky from the excess of lines over the null, drawn
again on the lines that hold in every exposure, and the per-line data, peaks and dips,
from which every number here follows.
\item The lines that recur at one sky position over several plate-nights, among them a
Ly$\alpha$ emitter at $z=2.893$ found by its recurrence alone.
\end{enumerate}

% =============================================================================
\section{Data and method}\label{sec:method}
% =============================================================================

\emph{Data.} DR17~\cite{Abdurrouf2022} holds, for each BOSS and eBOSS plate and night,
an spPlate file with the flux, inverse variance and sky-model spectrum of every fibre on
a common grid of $10^{-4}$ in $\log_{10}\lambda$, and the plug map that names each
fibre's target type. The fibres of type SKY give the sky spectra, whose rows we read from
the archive by byte range, one plate at a time: $4131$ files of the $3853$ plates.

\emph{Detection.} Each sky spectrum is scaled to its own robust noise: the flux times
the square root of the inverse variance, divided by $1.4826$ times its median absolute
deviation, with pixels the pipeline masks given zero weight. A filter matched to a
Gaussian of $\sigma=1.1$ pixels, the instrumental profile, gives the significance of a
line at each pixel, and a peak at $6\sigma$ or more is kept unless one of four tests sets
it aside: a masked pixel within one of it; an airglow line, where the plate's median
sky model exceeds three times its running continuum, widened by two pixels; two or more
other sky fibres of the same plate at $3\sigma$ within a pixel, which marks a residual of
the sky subtraction; or a width, from the peak's second moment, outside $0.6$ to $2$
times the profile. A peak that passes all four is clean. The same search run on the
negated spectra gives dips of the same shape. Each peak and dip that passes the first,
second and fourth tests is kept with the number of other fibres at $3\sigma$ beside it,
so the rule of one other fibre, stricter and less complete, is applied to the same
detections for comparison.

\emph{Tests after detection.} Three tests act on the clean lines of each sign alike. A
line that the same fibre holds within three pixels on another plate, at a sky position
more than $1.5''$ away, belongs to the detector; plates of one field can return a fibre
to the same position, and a line it holds there is the sky's. Two or more lines of one
fibre within $4\times10^{-4}(1+z)$ of the vacuum wavelengths of
[O\,\textsc{ii}]~$3727$, H$\beta$, [O\,\textsc{iii}]~$4959$ and $5007$, H$\alpha$,
[N\,\textsc{ii}]~$6583$, [S\,\textsc{ii}]~$6716$ and $6731$ or Ly$\alpha$ at one redshift
above $-0.002$, which admits Galactic nebulae, mark a galaxy, provided that one of them is
[O\,\textsc{ii}], H$\beta$, [O\,\textsc{iii}]~$5007$ or Ly$\alpha$, or H$\alpha$ beside
one of those, and that [O\,\textsc{iii}]~$4959$ appears only with $5007$ wherever $5007$
falls in the spectrum, at a flux ratio $5007/4959$ of $1.5$ to $5$ about the $2.98$ of the
transition probabilities. The dips, which no galaxy makes, measure how often residuals meet
this test by chance. A spectrum that holds clean lines of both signs outside the first test
is defective, and its lines of either sign are set aside. Each galaxy's position is
matched within $1.5''$ to the Legacy Surveys DR10 catalogue~\cite{Dey2019} and to DESI
DR1's redshifts~\cite{DESI2025} through the Astro Data Lab~\cite{Fitzpatrick2014}.

\emph{Counterparts.} A galaxy too faint in its other lines shows one line, and the
imaging may hold it. For every line at $8\sigma$ that the tests leave, peaks and dips
alike, the Legacy Surveys DR10 sources within $1.5''$ of the fibre and of the position
$30''$ north of it are read from the Astro Data Lab and the Legacy Survey viewer's
catalogue service; the excess at the fibres over the northern positions counts the lines
that come from catalogued sources.

\emph{Exposures.} An spPlate spectrum combines several exposures, each recorded in a
calibrated frame of the blue camera and one of the red. Every line at $8\sigma$ that the
tests leave, and every galaxy line at $8\sigma$, is measured again in each frame that
covers it, with a Gaussian of the frames' instrumental profile, $\sigma=1.2$ pixels,
placed within three pixels of the line's wavelength. A line that one frame carries, at
$5\sigma$ or more with over $80$ per cent of the summed squared significance, is a cosmic
ray or a detector event; one at $4\sigma$ in one camera and below $2\sigma$ in the other,
where the cameras overlap, is that camera's; one at $2.5\sigma$ or more in two or more
frames is persistent; the rest are unclear, and a line that fewer than two frames cover
is left unjudged.

\emph{Recurring positions.} A plate observed on several nights keeps its sky fibres where
they are, and later plates of a field return to some of its positions. Clean lines at one
position, friends of friends within $1.5''$, recur when they fall on two or more
plate-nights within $3\times10^{-4}$ in $\log\lambda$; a line that recurs at the same
wavelength at another position as well is the sky's. The dips, counted the same way, are
the null.

% =============================================================================
\section{Results}\label{sec:results}
% =============================================================================

\subsection{Peaks and dips}

The $381{,}943$ sky spectra hold $115{,}169$ emission-like peaks at $6\sigma$ and
$75{,}591$ absorption-like dips. The four tests leave $3217$ peaks and $4108$ dips, of
which $970$ and $1479$ reach $8\sigma$. Table~\ref{tab:counts} follows both signs through
the tests after detection, under the two-fibre rule and the one-fibre rule.

\begin{table}[t]
\centering\small
\caption{Peaks and dips at each stage, under the rule of two other fibres and of one.
Each later row applies its test to both signs; the counterpart rows give the lines with
a Legacy Surveys DR10 source within $1.5''$ at the fibre and at the position $30''$
north, among those the catalogue covers, and the persistent rows the lines of the
two-fibre rule that hold in two or more of the exposures' frames.}
\label{tab:counts}
\begin{tabular}{lrrrr}
\toprule
 & \multicolumn{2}{c}{Two other fibres} & \multicolumn{2}{c}{One other fibre} \\
 & Peaks & Dips & Peaks & Dips \\
\midrule
Clean at $6\sigma$ & $3217$ & $4108$ & $1633$ & $1874$ \\
\quad fixed to a fibre & $762$ & $1813$ & $242$ & $847$ \\
\quad at a galaxy's redshift & $369$ & $22$ & $144$ & $8$ \\
\quad with lines of both signs & $30$ & $38$ & $17$ & $14$ \\
\quad the rest & $2056$ & $2235$ & $1230$ & $1005$ \\
Clean at $8\sigma$ & $970$ & $1479$ & $521$ & $783$ \\
\quad fixed to a fibre & $287$ & $1114$ & $119$ & $537$ \\
\quad at a galaxy's redshift & $188$ & $19$ & $78$ & $6$ \\
\quad with lines of both signs & $10$ & $10$ & $9$ & $4$ \\
\quad the rest & $485$ & $336$ & $315$ & $236$ \\
\quad\quad source at the fibre & $181$ & $10$ & $138$ & $6$ \\
\quad\quad source $30''$ north & $21$ & $28$ & $12$ & $15$ \\
\quad\quad net of counterparts & $323$ & $354$ & $188$ & $245$ \\
\quad\quad persistent & $359$ & $241$ & & \\
\quad\quad\quad net of counterparts & $202$ & $253$ & & \\
\bottomrule
\end{tabular}
\end{table}

\subsection{Lines fixed to a fibre}

A line that one fibre holds on several plates at one wavelength and at different
positions belongs to the detector, not the sky: $762$ clean peaks, $287$ at $8\sigma$, in
$131$ fibres, and $1813$ dips, $1114$ at $8\sigma$, recur in the same fibre within three
pixels on other plates. Fibre $832$ near $9280$\,\AA\ carries $42$ of the peaks, on the red
camera where the airglow forest is densest, and fibres $10$ at $6600$\,\AA\ and $210$ at
$4460$\,\AA\ carry $19$ and $17$. The test within a plate cannot find these, since each
plate holds the fibre once, and three quarters of the dips at $8\sigma$ are of this kind.
The condition on position keeps on the sky $36$ peaks and one dip that a fibre holds where
a later plate of the field returns it to the same position.

\subsection{Galaxies}

The redshift test marks $156$ fibres, against $11.6\pm3.5$ when the peaks' wavelengths
are shuffled among the $325$ fibres that hold two or more, each fibre keeping its count,
over $200$ realisations, and $4$ among the dips. Without its requirements on H$\alpha$
and the [O\,\textsc{iii}] ratio it would mark $187$, and $30$ among the dips: H$\alpha$
with [N\,\textsc{ii}] or [S\,\textsc{ii}] alone ($16$ fibres among the peaks and $11$
among the dips), [O\,\textsc{iii}]~$4959$ without $5007$ ($11$ and $6$) and the
[O\,\textsc{iii}] doublet at a ratio outside $1.5$ to $5$ ($4$ and $9$), the forms the
residuals of the sky subtraction take. Some of the fibres set aside repeat features at
the same observed wavelengths on other plates, at $9842$ and $10097$\,\AA, $9842$ and
$9886$\,\AA\ and $9906$ and $10113$\,\AA\ among others, the residuals the subtraction
leaves between the strong OH lines of the red camera.

Two fibres hold lines at zero redshift, the nebular emission of the Galaxy where a fibre
crosses it. The other $154$ are galaxies (Table~\ref{tab:galaxies}) at $z=0.016$ to
$1.069$, median $0.441$ and quartiles $0.33$ and $0.64$, found by
[O\,\textsc{iii}]~$5007$ in $139$, [O\,\textsc{ii}] in $71$, H$\alpha$ in $56$,
[O\,\textsc{iii}]~$4959$ in $35$, H$\beta$ in $32$ and [S\,\textsc{ii}] and
[N\,\textsc{ii}] in one each. The $19$ found by the [O\,\textsc{iii}] doublet alone have
flux ratios $5007/4959$ of $1.8$ to $4.1$, median $2.7$, about the $2.98$ of the
transition probabilities. The twelve whose lines all lie above $8300$\,\AA, in the red
forest, are each H$\beta$ with [O\,\textsc{iii}] or the [O\,\textsc{iii}] doublet. The
[O\,\textsc{iii}]~$5007$ fluxes run from $3.0\times10^{-17}$ to
$9.8\times10^{-16}$\,\flux\ with a median of $8.5\times10^{-17}$. Six galaxies were
observed on two plates or nights, in different fibres, and give the same lines and
redshift each time.

One hundred and twenty-two of the $154$ have a Legacy Surveys DR10 source within
$1.5''$, of $r=21.8$ to $24.5$ and median $23.1$: $56$ fitted as round exponentials, $38$
as point sources, $24$ as exponentials, three as de Vaucouleurs profiles and one as a
S\'ersic profile. The $32$ without a source lie below the catalogue's depth, so their
lines are all the imaging and spectroscopy know of them. None has a DESI DR1 redshift
within $1.5''$: DESI's targets are chosen from the same imaging, and the sky fibres'
galaxies lie below the selection.

\subsection{What remains}

At $8\sigma$, $10$ peaks and $10$ dips lie in spectra that hold clean lines of both signs,
and $485$ peaks and $336$ dips are neither fixed to a fibre, at a galaxy's redshift nor in
such a spectrum. They are single lines in single fibres: a galaxy whose second line falls
on an airglow line, below threshold or beyond the spectrum, a residual of the sky
subtraction, or a transmitter would each appear so. Legacy Surveys sources lie within
$1.5''$ of $181$ of the $480$ peaks the catalogue covers against $21$ at positions $30''$
north, and of $10$ of the $330$ dips against $28$; scaled to every line, $162$ of the
peaks are single lines of galaxies the imaging catalogues, while the dips fall at
catalogued sources less often than at random positions. Net of the counterparts, $323$
peaks remain against the $336$ dips, or $354$ with the dips' shortfall at sources
restored, so the peaks fall short of the dips by $13$ to $31$, with an error of $33$.
Under the one-fibre rule the same counts give $188$ peaks against $236$ to $245$ dips.

The excess of peaks before the counterparts are counted lies at widths of $1.1$ to
$1.7$ pixels, the instrumental profile and a little wider, $210$ peaks against $101$
dips, where a galaxy's line falls; below $0.9$ pixels, where a cosmic ray the pipeline
leaves unmasked would fall, $10$ peaks lie against $13$ dips. The width test alone does
not set aside a single-pixel spike, since the second moment of the positive part of the
profile carries the noise beside it: in the package's end-to-end test a one-pixel spike
at $16\sigma$ on synthetic noise passes at a width of $0.82$ pixels. The dips' tail is
the longer, $47$ above $16\sigma$ against $27$ peaks, many of them where the blue and
red cameras' spectra are joined near $6000$\,\AA. Between $6$ and $8\sigma$ the dips
outnumber the peaks, $1899$ against $1571$, so the residuals do not make the two signs
equally at low significance, and the comparison is drawn at $8\sigma$, where they do to
within the counts.

\subsection{Exposure by exposure}

Read in the frames of their exposures, $179$ of the $188$ galaxy lines at $8\sigma$ are
persistent, seven unclear and two carried by one exposure or one camera, as lines on the
sky must be, while four of the $19$ dips at a galaxy's redshift are. Of the $485$ peaks
that remain, $359$ are persistent, $92$ unclear, $25$ carried by one exposure and one by
one camera, and eight lie on fewer than two frames; of the $336$ dips, $241$ are
persistent, $35$ unclear, $37$ carried by one exposure and $16$ by one camera, and seven
lie on fewer than two frames. Legacy Surveys sources lie within $1.5''$ of $175$ of the
$358$ persistent peaks the catalogue covers against $18$ at the control positions, and of
$7$ of the $236$ persistent dips against $19$, so net of the counterparts $202$ persistent
peaks stand against $253$ persistent dips, $52\pm29$ short.

The six strongest peaks that remain, at $36$ to $52\sigma$, have no Legacy Surveys source.
Four are carried by one exposure and one lies on a single frame. The sixth, at
$4961.4$\,\AA\ in fibre $562$ of plate $5294$, is [O\,\textsc{iii}]~$4959$ of the
planetary nebula Abell~31 (PN~G219.1$+$31.2), whose spectrum fills the fibre from
[Ne\,\textsc{iii}]~$3869$ to [S\,\textsc{iii}]~$9531$ at zero redshift: its brighter lines
reach $3\sigma$ in two or more of the plate's other sky fibres and are set aside as shared,
the sky fibres off the nebula hold them as dips where the sky model took in its light, and
five fainter lines, H$\gamma$, [O\,\textsc{iii}]~$4959$, [Ar\,\textsc{iii}]~$7136$ and
[S\,\textsc{iii}]~$9069$ and $9531$, remain among the peaks.

\subsection{Lines that recur at one position}

Clean lines fall on two or more plate-nights at one position at $53$ positions among the
peaks and $28$ among the dips. One line recurs at $37$ of the peak positions and at three
of the dip positions, and the recurring peaks are $70$ of the $485$ at $8\sigma$, against
three dips. Five of the $37$, all between $7204$ and $7218$\,\AA, recur at another position
as well and are the sky's.

The line found most often lies at RA $213.5584^\circ$, Dec $+53.6108^\circ$ in the field
of the SDSS reverberation-mapping project~\cite{Shen2015}, where plates $7338$, $7339$ and
$7340$, observed on many nights, and three others put a sky fibre in $94$ spectra between
2013 and 2020. A line at $4732.6$\,\AA\ is detected in $81$ of them, clean in $54$ and
among the peaks that remain at $8\sigma$ in $41$, with no Legacy Surveys source at the
fibre. The $94$ spectra stacked by inverse variance hold it at
$52\sigma$, centred at $4732.7$\,\AA\ with a flux of $8.5\times10^{-17}$\,\flux\ and a
width of $287$\,km\,s$^{-1}$ against the $179$\,km\,s$^{-1}$ of the instrumental profile,
and no other line. Read as [O\,\textsc{ii}] at $z=0.269$, it puts H$\beta$,
[O\,\textsc{iii}]~$5007$ and H$\alpha$ at $6172$, $6357$ and $8333$\,\AA, where the stack
holds $-0.02$, $0.00$ and $0.00$ of its flux with errors of $0.01$, while a star-forming
galaxy's H$\alpha$ is as bright as its [O\,\textsc{ii}] or brighter. Read as
C\,\textsc{iv} at $z=2.054$, it puts He\,\textsc{ii} and C\,\textsc{iii}] where the stack
holds $-0.01$ and $0.01\pm0.01$, and read as Mg\,\textsc{ii} at $z=0.690$, it puts
[O\,\textsc{ii}] and [O\,\textsc{iii}] where the stack holds $-0.06\pm0.03$ and
$0.05\pm0.02$; an active nucleus would show those lines, and its Mg\,\textsc{ii} is
broad. Read as Ly$\alpha$ at $z=2.893$, it puts N\,\textsc{v}, C\,\textsc{iv},
He\,\textsc{ii} and C\,\textsc{iii}] where the stack holds $0.00$ to $0.01\pm0.01$, lines
that star-forming galaxies at $z\approx3$ show at a few per cent of
Ly$\alpha$~\cite{Shapley2003}, and its intrinsic width, about $220$\,km\,s$^{-1}$, is that
of a Ly$\alpha$ emitter: a galaxy at $z=2.893$ below the imaging's depth, resolved where a
transmitter's line is not.

The second, [O\,\textsc{iii}]~$5007$ at $7775.0$\,\AA, recurs at RA $29.4755^\circ$, Dec
$+27.5883^\circ$ on all seven plate-nights of plates $7691$ and $10000$ that hold a clean
line there, with H$\beta$ beside it on one, which places it in
Table~\ref{tab:galaxies} at $z=0.5525$. The other $35$ positions repeat a line on two or
three plate-nights, and $15$ of their $26$ lines at $8\sigma$ have a Legacy Surveys source
within $1.5''$ against one at the control positions.

\subsection{Sensitivity}

The line flux at the $6\sigma$ threshold, read from each clean peak's flux and
significance, has a median of $6.0\times10^{-17}$\,\flux\ and a $10$--$90$ per cent
range of $3.4\times10^{-17}$ to $2.7\times10^{-16}$: $6.2\times10^{-17}$ from $3600$ to
$5000$\,\AA, $5.0\times10^{-17}$ from $5000$ to $7000$, $4.8\times10^{-17}$ from $7000$
to $9000$ and $1.5\times10^{-16}$ beyond $9000$\,\AA, where the airglow lines crowd. At
$8\sigma$ the median is $8.0\times10^{-17}$\,\flux. An isotropic source giving that flux
has a luminosity of $9.5\times10^{16}$\,W at $10$\,pc, $9.5\times10^{18}$\,W at
$100$\,pc and $9.5\times10^{20}$\,W at $1$\,kpc, and a transmitter aimed at Earth through
a diffraction-limited $10$\,m aperture at the median wavelength of $7013$\,\AA\ needs
$37$\,W, $3.7$\,kW and $370$\,kW; its beam is $0.14$\,au wide at $10$\,pc and $1.4$\,au
at $100$\,pc, so within $100$\,pc it must be held on the Earth's position at arrival,
and beyond it a beam on the Sun covers the Earth as well.

Lines of the instrumental profile injected into the sky fibres of $40$ plates spread
through the list, one per fibre at a random unmasked pixel off the airglow lines and
$1727$ at each strength, are recovered as clean peaks under the two-fibre rule in $56$
per cent of cases at the threshold and $70$ to $74$ per cent at $1.5$ to $5$ times it,
and under the one-fibre rule in $30$ and $36$ to $38$ per cent; the loss is the test
against the other fibres, whose residuals reach $3\sigma$ within a pixel of a given
position in some fibre often enough to read a line as shared. On synthetic plates of
white noise the same detector recovers lines at $10\sigma$ in $99$ and $84$ per cent of
cases under the two rules, so the loss in the data is the residuals'. At $8\sigma$, one
third above the threshold, the two-fibre rule recovers about $0.68$ of lines.

The fibres are $2''$ across and cover $0.093$\,deg$^2$ in all, so the search bounds the
density of transmitters on the sky rather than any chosen star. The excess of peaks over
dips at $8\sigma$, below zero with an error of $33$, puts the number of transmitter lines
below $54$ at $95$ per cent confidence, the larger of zero and the excess plus $1.645$
times its error, and with $0.68$ of lines recovered, fewer than about $850$ transmitters
per square degree give lines above $8\times10^{-17}$\,\flux\ at these wavelengths,
$3.5\times10^{7}$ over the sky. A transmitter that shines through a plate's exposures
persists in its frames as $95$ per cent of the galaxies' lines do, which leaves $0.65$ of
its lines recovered at $8\sigma$, and the persistent peaks, $52\pm29$ short of the
persistent dips, put fewer than $47$ such lines in the data and fewer than about $790$
such transmitters on each square degree.

\begin{center}\small
\setlength{\tabcolsep}{4pt}
\begin{longtable}{lrrr>{\raggedright\arraybackslash}p{3.9cm}rl}
\caption{The $154$ galaxies and two Galactic fibres: plate, MJD and fibre; the fibre's
position; the redshift; the lines found at it; the strongest line's signal-to-noise
ratio; and the Legacy Surveys DR10 source within $1.5''$, with its profile type and
$r$.}\label{tab:galaxies}\\
\toprule
Plate-MJD-fibre & RA & Dec & $z$ & Lines & S/N & Legacy Surveys \\
\midrule
\endfirsthead
\toprule
Plate-MJD-fibre & RA & Dec & $z$ & Lines & S/N & Legacy Surveys \\
\midrule
\endhead
\bottomrule
\endfoot
8169-57071-296 & 127.14855 & +57.12516 & -0.0002 & [N\,\textsc{ii}], [O\,\textsc{iii}]\,5007, [S\,\textsc{ii}]\,6731 & 7.5 & --\\
6602-56571-766 & 24.02404 & +34.10890 & 0.0001 & [O\,\textsc{ii}], [S\,\textsc{ii}]\,6716 & 8.2 & outside\\
11288-58401-543 & 351.36997 & +6.09305 & 0.0164 & H$\alpha$, [O\,\textsc{iii}]\,5007 & 8.7 & --\\
4020-55332-26 & 226.04718 & -2.39425 & 0.0954 & H$\alpha$, [O\,\textsc{iii}]\,5007 & 13.5 & --\\
4563-55617-722 & 154.88368 & +38.46850 & 0.1085 & [O\,\textsc{iii}]\,4959, [O\,\textsc{iii}]\,5007 & 14.9 & --\\
8175-57039-450 & 150.61150 & +50.64228 & 0.1416 & H$\alpha$, [O\,\textsc{iii}]\,5007 & 10.8 & round exp., $r=22.5$\\
7512-56777-45 & 134.07495 & +48.20114 & 0.1633 & H$\alpha$, [O\,\textsc{iii}]\,5007 & 13.6 & --\\
3877-55365-396 & 225.00402 & +28.04808 & 0.1817 & H$\beta$, [O\,\textsc{iii}]\,5007 & 10.4 & round exp., $r=23.3$\\
6013-56074-554 & 213.58866 & +24.33099 & 0.1895 & H$\alpha$, [O\,\textsc{iii}]\,5007 & 7.5 & point, $r=23.6$\\
4057-55357-108 & 245.33337 & +20.30591 & 0.1997 & H$\beta$, [O\,\textsc{iii}]\,5007 & 10.5 & exp., $r=22.5$\\
9435-58108-594 & 17.53466 & +4.66545 & 0.1999 & H$\alpha$, [O\,\textsc{iii}]\,5007, [O\,\textsc{ii}] & 10.8 & round exp., $r=22.3$\\
4279-55508-88 & 357.85425 & +2.95924 & 0.2248 & H$\alpha$, [O\,\textsc{iii}]\,5007 & 16.3 & round exp., $r=22.2$\\
7690-57046-103 & 28.95435 & +25.09805 & 0.2249 & H$\alpha$, [O\,\textsc{iii}]\,5007 & 13.1 & round exp., $r=22.3$\\
5378-56011-408 & 174.56869 & +14.12900 & 0.2277 & [O\,\textsc{iii}]\,4959, [O\,\textsc{iii}]\,5007 & 19.4 & --\\
7822-57041-354 & 41.57260 & -1.06696 & 0.2313 & H$\alpha$, [O\,\textsc{ii}] & 7.2 & exp., $r=22.0$\\
6793-56430-248 & 242.35395 & +54.12879 & 0.2351 & H$\alpha$, [O\,\textsc{iii}]\,5007 & 9.5 & exp., $r=22.5$\\
5404-56013-406 & 187.52239 & +14.68962 & 0.2357 & H$\alpha$, [O\,\textsc{iii}]\,5007 & 11.3 & round exp., $r=24.0$\\
11061-58428-6 & 26.43771 & +11.21563 & 0.2360 & H$\alpha$, [O\,\textsc{iii}]\,5007 & 7.7 & round exp., $r=22.5$\\
5873-56035-868 & 158.02560 & +22.15247 & 0.2445 & H$\alpha$, [O\,\textsc{iii}]\,5007, [O\,\textsc{ii}] & 8.9 & round exp., $r=23.4$\\
6689-56396-676 & 177.85168 & +54.09125 & 0.2483 & H$\alpha$, [O\,\textsc{iii}]\,5007 & 7.8 & round exp., $r=22.3$\\
11352-58456-248 & 156.58603 & +27.53902 & 0.2508 & H$\alpha$, [O\,\textsc{iii}]\,5007 & 12.6 & round exp., $r=23.3$\\
7311-57038-991 & 132.07073 & +50.44043 & 0.2680 & H$\alpha$, [O\,\textsc{iii}]\,5007 & 12.8 & --\\
7512-56777-675 & 132.07073 & +50.44043 & 0.2680 & H$\alpha$, [O\,\textsc{iii}]\,5007 & 13.3 & --\\
7699-57574-95 & 355.68318 & +25.99462 & 0.2812 & H$\alpha$, [O\,\textsc{iii}]\,5007 & 14.1 & exp., $r=22.6$\\
10901-58397-179 & 255.53170 & +32.14931 & 0.2857 & H$\alpha$, [O\,\textsc{iii}]\,5007 & 9.7 & exp., $r=22.6$\\
7835-56986-322 & 33.83880 & -3.32220 & 0.2939 & [O\,\textsc{iii}]\,4959, [O\,\textsc{iii}]\,5007 & 14.6 & round exp., $r=23.3$\\
7390-57130-18 & 163.85952 & +47.34356 & 0.2948 & [O\,\textsc{iii}]\,5007, [O\,\textsc{ii}] & 13.2 & round exp., $r=22.8$\\
5046-56243-614 & 343.65339 & +13.46747 & 0.2988 & [O\,\textsc{iii}]\,5007, [O\,\textsc{ii}] & 8.0 & round exp., $r=22.7$\\
5002-55710-284 & 261.59575 & +30.59913 & 0.3008 & H$\alpha$, [O\,\textsc{iii}]\,5007 & 10.1 & round exp., $r=22.3$\\
9575-57786-786 & 153.75520 & +17.37566 & 0.3050 & H$\alpha$, H$\beta$, [O\,\textsc{iii}]\,4959, [O\,\textsc{iii}]\,5007 & 16.6 & --\\
11561-58485-827 & 2.04539 & +11.85705 & 0.3056 & H$\alpha$, [O\,\textsc{iii}]\,5007 & 9.0 & point, $r=23.1$\\
4040-55605-420 & 207.18925 & -0.83451 & 0.3062 & H$\alpha$, [O\,\textsc{iii}]\,5007 & 7.3 & round exp., $r=23.5$\\
5141-55746-113 & 318.86697 & +0.91202 & 0.3092 & H$\alpha$, [O\,\textsc{iii}]\,5007 & 9.2 & exp., $r=22.5$\\
7057-56593-810 & 41.11542 & -7.41403 & 0.3110 & H$\alpha$, H$\beta$, [O\,\textsc{iii}]\,4959, [O\,\textsc{iii}]\,5007, [O\,\textsc{ii}] & 19.0 & S\'ersic, $r=22.1$\\
4740-55651-178 & 174.48659 & +1.11383 & 0.3113 & H$\alpha$, H$\beta$, [S\,\textsc{ii}]\,6716 & 18.6 & round exp., $r=21.8$\\
7900-57358-798 & 4.13084 & -3.05673 & 0.3125 & H$\alpha$, [O\,\textsc{iii}]\,5007 & 12.5 & point, $r=23.5$\\
10440-58138-745 & 141.43832 & +31.17449 & 0.3134 & [O\,\textsc{iii}]\,4959, [O\,\textsc{iii}]\,5007 & 20.4 & point, $r=23.6$\\
9141-57663-653 & 321.07987 & -0.46193 & 0.3201 & H$\beta$, [O\,\textsc{iii}]\,4959, [O\,\textsc{iii}]\,5007, [O\,\textsc{ii}] & 18.2 & round exp., $r=22.6$\\
6743-56385-218 & 202.05611 & +47.97118 & 0.3268 & H$\alpha$, [O\,\textsc{iii}]\,4959, [O\,\textsc{iii}]\,5007 & 16.4 & point, $r=24.1$\\
11109-58523-950 & 124.27555 & +22.24246 & 0.3279 & H$\alpha$, [O\,\textsc{iii}]\,5007 & 9.0 & round exp., $r=22.2$\\
4174-55659-442 & 255.95617 & +21.05868 & 0.3289 & H$\alpha$, [O\,\textsc{iii}]\,5007 & 7.8 & round exp., $r=23.5$\\
5198-55823-408 & 252.65262 & +35.90116 & 0.3395 & [O\,\textsc{iii}]\,5007, [O\,\textsc{ii}] & 13.2 & round exp., $r=23.9$\\
8370-57520-708 & 175.74678 & +44.52171 & 0.3412 & H$\alpha$, [O\,\textsc{iii}]\,5007 & 12.1 & round exp., $r=23.4$\\
8370-57732-718 & 175.74678 & +44.52171 & 0.3412 & H$\alpha$, [O\,\textsc{iii}]\,5007 & 10.1 & round exp., $r=23.4$\\
10738-58251-439 & 240.90571 & +35.18927 & 0.3416 & H$\alpha$, [O\,\textsc{ii}] & 12.0 & round exp., $r=22.7$\\
7095-56625-482 & 160.11325 & +58.20823 & 0.3498 & [O\,\textsc{iii}]\,4959, [O\,\textsc{iii}]\,5007, [O\,\textsc{ii}] & 23.4 & round exp., $r=22.6$\\
6587-56537-574 & 349.59914 & +27.54950 & 0.3516 & [O\,\textsc{iii}]\,5007, [O\,\textsc{ii}] & 9.0 & exp., $r=22.1$\\
10917-58252-746 & 241.41881 & +32.47987 & 0.3524 & H$\alpha$, [O\,\textsc{iii}]\,5007 & 8.5 & round exp., $r=23.5$\\
6633-56369-490 & 183.16471 & +41.81318 & 0.3568 & [O\,\textsc{iii}]\,4959, [O\,\textsc{iii}]\,5007 & 15.3 & round exp., $r=23.6$\\
4543-55888-240 & 11.54328 & +6.28194 & 0.3584 & [O\,\textsc{iii}]\,5007, [O\,\textsc{ii}] & 7.6 & point, $r=23.2$\\
9382-57746-452 & 34.84357 & +0.03436 & 0.3624 & [O\,\textsc{iii}]\,5007, [O\,\textsc{ii}] & 14.8 & de Vauc., $r=22.2$\\
6586-56485-366 & 336.95291 & +28.14091 & 0.3655 & [O\,\textsc{iii}]\,5007, [O\,\textsc{ii}] & 9.8 & round exp., $r=21.9$\\
6289-56559-308 & 348.17625 & +27.40934 & 0.3717 & [N\,\textsc{ii}], [O\,\textsc{ii}] & 8.4 & point, $r=22.7$\\
10748-58224-776 & 223.64392 & +38.18774 & 0.3775 & [O\,\textsc{iii}]\,5007, [O\,\textsc{ii}] & 7.7 & round exp., $r=23.4$\\
4643-55946-153 & 138.35170 & +37.83216 & 0.3833 & H$\alpha$, [O\,\textsc{iii}]\,5007 & 9.5 & round exp., $r=22.8$\\
4016-55632-414 & 225.44444 & -0.61182 & 0.3888 & H$\alpha$, [O\,\textsc{iii}]\,5007 & 8.6 & --\\
9442-58076-984 & 9.95057 & +2.65555 & 0.3919 & H$\alpha$, [O\,\textsc{iii}]\,4959, [O\,\textsc{iii}]\,5007 & 14.7 & --\\
7701-57356-625 & 351.51853 & +27.44095 & 0.3932 & H$\alpha$, [O\,\textsc{iii}]\,5007 & 10.2 & --\\
11059-58515-987 & 23.81695 & +17.02814 & 0.3960 & H$\alpha$, [O\,\textsc{ii}] & 11.2 & round exp., $r=22.8$\\
9162-58040-242 & 322.00060 & +0.38206 & 0.3972 & H$\alpha$, [O\,\textsc{iii}]\,5007 & 9.5 & point, $r=23.1$\\
7631-57279-635 & 18.95434 & +22.18753 & 0.3977 & H$\alpha$, [O\,\textsc{iii}]\,5007 & 10.2 & --\\
8526-57898-317 & 244.45013 & +41.51516 & 0.3982 & H$\alpha$, H$\beta$, [O\,\textsc{iii}]\,4959, [O\,\textsc{iii}]\,5007 & 14.5 & point, $r=24.0$\\
6044-56090-107 & 244.45014 & +41.51515 & 0.3983 & H$\alpha$, [O\,\textsc{iii}]\,4959, [O\,\textsc{iii}]\,5007 & 17.6 & point, $r=24.0$\\
8792-57364-446 & 23.51387 & -1.05226 & 0.4020 & H$\alpha$, [O\,\textsc{iii}]\,5007 & 7.0 & round exp., $r=23.6$\\
4656-55828-488 & 28.11478 & +11.68275 & 0.4028 & H$\alpha$, [O\,\textsc{iii}]\,5007 & 8.4 & point, $r=23.7$\\
5694-56213-351 & 23.77663 & +25.07137 & 0.4034 & H$\alpha$, [O\,\textsc{ii}] & 7.7 & --\\
3878-55361-470 & 221.72976 & +28.42111 & 0.4085 & H$\alpha$, [O\,\textsc{iii}]\,5007 & 8.7 & round exp., $r=23.4$\\
8211-57423-986 & 205.78650 & +50.99986 & 0.4098 & H$\alpha$, [O\,\textsc{iii}]\,5007 & 12.0 & --\\
6423-56313-38 & 178.98923 & +22.51120 & 0.4127 & H$\alpha$, [O\,\textsc{iii}]\,5007 & 10.7 & round exp., $r=23.1$\\
4038-55363-444 & 210.29613 & -3.23692 & 0.4129 & H$\alpha$, [O\,\textsc{iii}]\,5007 & 14.7 & exp., $r=22.6$\\
5777-56280-948 & 140.43369 & +26.23998 & 0.4152 & H$\alpha$, [O\,\textsc{iii}]\,5007 & 7.0 & round exp., $r=22.7$\\
8750-58401-690 & 10.33064 & +6.25421 & 0.4302 & [O\,\textsc{iii}]\,5007, [O\,\textsc{ii}] & 15.2 & point, $r=23.3$\\
9554-57807-188 & 140.49693 & +18.22247 & 0.4315 & H$\alpha$, [O\,\textsc{iii}]\,5007 & 8.5 & point, $r=23.6$\\
5307-55998-786 & 140.49693 & +18.22247 & 0.4315 & H$\alpha$, [O\,\textsc{iii}]\,5007 & 10.2 & point, $r=23.6$\\
6402-56334-335 & 178.75863 & +27.52210 & 0.4372 & [O\,\textsc{iii}]\,4959, [O\,\textsc{iii}]\,5007 & 12.4 & point, $r=23.9$\\
6588-56536-822 & 344.67859 & +25.91339 & 0.4374 & [O\,\textsc{iii}]\,5007, [O\,\textsc{ii}] & 9.4 & round exp., $r=23.2$\\
8732-57418-829 & 35.93424 & +5.21363 & 0.4385 & [O\,\textsc{iii}]\,4959, [O\,\textsc{iii}]\,5007 & 20.8 & point, $r=23.6$\\
7114-56748-520 & 169.81510 & +65.92926 & 0.4398 & [O\,\textsc{iii}]\,5007, [O\,\textsc{ii}] & 11.6 & round exp., $r=23.0$\\
11554-58456-170 & 353.86311 & +9.93507 & 0.4408 & [O\,\textsc{iii}]\,5007, [O\,\textsc{ii}] & 9.9 & point, $r=23.7$\\
6165-56208-788 & 353.86311 & +9.93508 & 0.4408 & [O\,\textsc{iii}]\,4959, [O\,\textsc{iii}]\,5007 & 16.5 & point, $r=23.7$\\
4641-55947-314 & 138.90323 & +40.53916 & 0.4416 & [O\,\textsc{iii}]\,5007, [O\,\textsc{ii}] & 12.2 & exp., $r=22.4$\\
7646-57570-208 & 332.07143 & +27.39413 & 0.4485 & [O\,\textsc{iii}]\,4959, [O\,\textsc{iii}]\,5007, [O\,\textsc{ii}] & 15.1 & point, $r=23.8$\\
6487-56362-880 & 198.02095 & +31.92434 & 0.4487 & [O\,\textsc{iii}]\,4959, [O\,\textsc{iii}]\,5007, [O\,\textsc{ii}] & 12.7 & point, $r=23.3$\\
10914-58257-755 & 248.98563 & +33.41463 & 0.4504 & [O\,\textsc{iii}]\,5007, [O\,\textsc{ii}] & 12.7 & round exp., $r=22.4$\\
11380-58462-704 & 150.33497 & +31.84991 & 0.4564 & [O\,\textsc{iii}]\,5007, [O\,\textsc{ii}] & 7.5 & round exp., $r=23.2$\\
7391-56781-675 & 160.40951 & +46.60769 & 0.4672 & H$\alpha$, [O\,\textsc{iii}]\,5007 & 13.1 & round exp., $r=22.9$\\
6516-56571-602 & 355.05918 & +27.08038 & 0.4683 & H$\beta$, [O\,\textsc{iii}]\,5007 & 13.4 & round exp., $r=23.4$\\
8497-58161-40 & 221.93867 & +42.20213 & 0.4761 & H$\beta$, [O\,\textsc{iii}]\,5007 & 31.3 & --\\
6203-56266-288 & 12.00438 & +14.31752 & 0.4827 & [O\,\textsc{iii}]\,5007, [O\,\textsc{ii}] & 6.9 & round exp., $r=22.0$\\
6699-56411-882 & 169.24331 & +52.16521 & 0.4856 & H$\alpha$, [O\,\textsc{iii}]\,5007, [O\,\textsc{ii}] & 8.6 & point, $r=23.0$\\
9610-58137-172 & 133.00640 & +27.21546 & 0.4895 & [O\,\textsc{iii}]\,5007, [O\,\textsc{ii}] & 7.0 & point, $r=22.8$\\
6634-56367-520 & 179.94605 & +41.95528 & 0.4895 & H$\alpha$, [O\,\textsc{iii}]\,5007 & 8.4 & point, $r=23.0$\\
6484-56358-594 & 193.85918 & +31.69102 & 0.4952 & [O\,\textsc{iii}]\,4959, [O\,\textsc{iii}]\,5007 & 14.0 & --\\
11071-58429-332 & 24.19181 & +6.40044 & 0.5064 & [O\,\textsc{iii}]\,5007, [O\,\textsc{ii}] & 9.2 & round exp., $r=23.1$\\
8279-57063-174 & 127.02745 & +43.19890 & 0.5109 & [O\,\textsc{iii}]\,5007, [O\,\textsc{ii}] & 7.3 & round exp., $r=23.3$\\
7107-56740-796 & 174.80896 & +62.15769 & 0.5146 & [O\,\textsc{iii}]\,5007, [O\,\textsc{ii}] & 9.1 & point, $r=23.1$\\
7606-56977-311 & 350.26431 & +21.01340 & 0.5177 & H$\beta$, [O\,\textsc{ii}] & 13.6 & exp., $r=22.1$\\
7860-57006-380 & 22.15587 & +1.18941 & 0.5261 & H$\beta$, [O\,\textsc{iii}]\,5007 & 6.8 & exp., $r=21.9$\\
8204-58146-416 & 198.21603 & +55.42233 & 0.5277 & [O\,\textsc{iii}]\,5007, [O\,\textsc{ii}] & 7.7 & exp., $r=22.0$\\
4994-55739-370 & 258.61404 & +33.83769 & 0.5293 & [O\,\textsc{iii}]\,4959, [O\,\textsc{iii}]\,5007 & 18.7 & --\\
6422-56328-90 & 176.49067 & +21.57333 & 0.5318 & H$\beta$, [O\,\textsc{iii}]\,4959, [O\,\textsc{iii}]\,5007, [O\,\textsc{ii}] & 32.5 & point, $r=23.4$\\
3847-55298-412 & 185.87047 & -0.69507 & 0.5432 & H$\beta$, [O\,\textsc{iii}]\,4959, [O\,\textsc{iii}]\,5007 & 21.5 & --\\
3847-55588-418 & 185.87047 & -0.69507 & 0.5432 & H$\beta$, [O\,\textsc{iii}]\,4959, [O\,\textsc{iii}]\,5007 & 20.6 & --\\
3868-55360-202 & 219.81346 & +30.49225 & 0.5504 & H$\beta$, [O\,\textsc{ii}] & 11.9 & exp., $r=22.3$\\
7691-58487-687 & 29.47546 & +27.58831 & 0.5525 & H$\beta$, [O\,\textsc{iii}]\,5007 & 13.1 & --\\
9144-57666-648 & 328.28503 & +0.02365 & 0.5563 & H$\beta$, [O\,\textsc{ii}] & 14.4 & round exp., $r=22.0$\\
5313-55973-558 & 142.25359 & +13.82727 & 0.5702 & [O\,\textsc{iii}]\,5007, [O\,\textsc{ii}] & 8.7 & exp., $r=22.7$\\
4747-55652-835 & 180.80191 & +4.98357 & 0.5885 & [O\,\textsc{iii}]\,5007, [O\,\textsc{ii}] & 12.0 & exp., $r=22.3$\\
10752-58488-557 & 217.96798 & +39.67842 & 0.5980 & H$\beta$, [O\,\textsc{iii}]\,5007 & 6.9 & --\\
3947-55332-324 & 235.38928 & +23.87738 & 0.6051 & [O\,\textsc{iii}]\,5007, [O\,\textsc{ii}] & 8.3 & exp., $r=22.3$\\
11337-58431-190 & 330.43338 & +5.56656 & 0.6088 & [O\,\textsc{iii}]\,5007, [O\,\textsc{ii}] & 6.9 & exp., $r=22.7$\\
3940-55327-364 & 239.43312 & +23.99639 & 0.6203 & H$\beta$, [O\,\textsc{iii}]\,4959, [O\,\textsc{iii}]\,5007 & 28.3 & exp., $r=22.6$\\
4533-55855-414 & 2.28310 & +7.12388 & 0.6208 & [O\,\textsc{iii}]\,4959, [O\,\textsc{iii}]\,5007, [O\,\textsc{ii}] & 20.1 & round exp., $r=22.8$\\
8755-58138-198 & 25.24367 & +5.87333 & 0.6300 & [O\,\textsc{iii}]\,5007, [O\,\textsc{ii}] & 8.2 & round exp., $r=22.5$\\
6155-56245-438 & 347.48707 & +10.35645 & 0.6312 & [O\,\textsc{iii}]\,4959, [O\,\textsc{iii}]\,5007 & 13.7 & --\\
5191-56065-388 & 245.24949 & +34.60598 & 0.6395 & [O\,\textsc{iii}]\,4959, [O\,\textsc{iii}]\,5007 & 11.8 & --\\
9582-57806-566 & 143.08944 & +17.21156 & 0.6396 & [O\,\textsc{iii}]\,5007, [O\,\textsc{ii}] & 9.1 & de Vauc., $r=22.4$\\
3781-55243-146 & 141.78061 & -2.11227 & 0.6453 & [O\,\textsc{iii}]\,5007, [O\,\textsc{ii}] & 11.7 & point, $r=23.1$\\
9162-58040-928 & 322.94891 & +1.91979 & 0.6476 & [O\,\textsc{iii}]\,4959, [O\,\textsc{iii}]\,5007 & 16.9 & point, $r=23.9$\\
4779-55656-64 & 221.47363 & +3.87355 & 0.6515 & [O\,\textsc{iii}]\,5007, [O\,\textsc{ii}] & 8.0 & exp., $r=22.5$\\
4753-55674-152 & 188.64795 & +3.56641 & 0.6634 & [O\,\textsc{iii}]\,5007, [O\,\textsc{ii}] & 12.4 & round exp., $r=22.9$\\
10429-58145-627 & 143.02481 & +32.28740 & 0.6823 & [O\,\textsc{iii}]\,5007, [O\,\textsc{ii}] & 13.1 & point, $r=23.0$\\
6497-56329-300 & 213.29485 & +29.07750 & 0.6937 & H$\beta$, [O\,\textsc{iii}]\,5007, [O\,\textsc{ii}] & 27.6 & exp., $r=22.5$\\
4858-55686-359 & 221.21223 & +6.63193 & 0.7023 & [O\,\textsc{iii}]\,5007, [O\,\textsc{ii}] & 8.5 & exp., $r=23.1$\\
8388-57488-346 & 198.47179 & +41.25348 & 0.7037 & [O\,\textsc{iii}]\,5007, [O\,\textsc{ii}] & 7.3 & round exp., $r=23.0$\\
7747-58398-852 & 3.14290 & +35.45403 & 0.7039 & [O\,\textsc{iii}]\,5007, [O\,\textsc{ii}] & 7.7 & outside\\
4221-55443-510 & 7.27573 & +0.01951 & 0.7185 & [O\,\textsc{iii}]\,5007, [O\,\textsc{ii}] & 10.5 & exp., $r=22.9$\\
5739-55976-150 & 149.61539 & +48.66389 & 0.7191 & [O\,\textsc{iii}]\,4959, [O\,\textsc{iii}]\,5007 & 14.0 & --\\
7183-56658-450 & 25.90411 & -10.25796 & 0.7360 & H$\beta$, [O\,\textsc{iii}]\,4959, [O\,\textsc{iii}]\,5007 & 21.7 & exp., $r=22.7$\\
8234-57450-582 & 177.12140 & +58.45203 & 0.7437 & H$\beta$, [O\,\textsc{iii}]\,5007 & 8.2 & round exp., $r=22.9$\\
5863-56042-816 & 197.31227 & +18.68198 & 0.7451 & H$\beta$, [O\,\textsc{iii}]\,5007 & 8.7 & round exp., $r=22.5$\\
5130-55835-42 & 26.45194 & +14.74649 & 0.7455 & H$\beta$, [O\,\textsc{iii}]\,5007 & 19.9 & point, $r=23.6$\\
4348-55559-700 & 27.49099 & -2.31464 & 0.7476 & [O\,\textsc{iii}]\,5007, [O\,\textsc{ii}] & 30.7 & round exp., $r=23.3$\\
8763-57309-452 & 348.32325 & -0.06942 & 0.7518 & H$\beta$, [O\,\textsc{ii}] & 6.7 & --\\
5733-56575-654 & 138.00488 & +48.77073 & 0.7541 & H$\beta$, [O\,\textsc{ii}] & 8.0 & --\\
6881-56540-554 & 2.46761 & +23.78729 & 0.7575 & [O\,\textsc{iii}]\,4959, [O\,\textsc{iii}]\,5007 & 16.6 & point, $r=23.4$\\
6825-56717-460 & 202.51184 & +66.02136 & 0.7933 & H$\beta$, [O\,\textsc{ii}] & 9.1 & round exp., $r=23.2$\\
5979-56329-226 & 185.96476 & +23.12713 & 0.7962 & H$\beta$, [O\,\textsc{ii}] & 8.9 & point, $r=23.6$\\
5800-56279-368 & 149.00528 & +29.49835 & 0.8026 & [O\,\textsc{iii}]\,5007, [O\,\textsc{ii}] & 6.2 & round exp., $r=23.6$\\
7854-56989-316 & 5.58985 & +1.51164 & 0.8032 & [O\,\textsc{iii}]\,5007, [O\,\textsc{ii}] & 7.6 & de Vauc., $r=23.1$\\
4376-55863-124 & 330.24601 & -3.46081 & 0.8089 & [O\,\textsc{iii}]\,4959, [O\,\textsc{iii}]\,5007 & 25.1 & --\\
7142-56567-606 & 338.98836 & +33.05879 & 0.8099 & [O\,\textsc{iii}]\,4959, [O\,\textsc{iii}]\,5007 & 19.7 & point, $r=24.0$\\
4207-55475-812 & 345.33907 & +0.36194 & 0.8357 & [O\,\textsc{iii}]\,5007, [O\,\textsc{ii}] & 11.0 & point, $r=24.5$\\
5413-55976-488 & 191.68889 & +7.49739 & 0.8365 & [O\,\textsc{iii}]\,5007, [O\,\textsc{ii}] & 7.6 & round exp., $r=23.2$\\
10757-58223-520 & 215.56647 & +35.41421 & 0.8369 & [O\,\textsc{iii}]\,5007, [O\,\textsc{ii}] & 6.8 & point, $r=23.8$\\
4478-55600-376 & 126.61986 & +21.27518 & 0.8550 & [O\,\textsc{iii}]\,5007, [O\,\textsc{ii}] & 8.1 & point, $r=23.7$\\
5410-56016-644 & 191.05282 & +14.52098 & 0.8552 & [O\,\textsc{iii}]\,5007, [O\,\textsc{ii}] & 6.8 & point, $r=24.0$\\
5447-55958-124 & 210.25983 & +7.91491 & 0.8552 & [O\,\textsc{iii}]\,5007, [O\,\textsc{ii}] & 6.3 & round exp., $r=23.4$\\
4447-55542-449 & 122.63066 & +29.97605 & 0.8795 & [O\,\textsc{iii}]\,5007, [O\,\textsc{ii}] & 8.5 & point, $r=24.2$\\
5703-56190-288 & 14.54918 & +19.37136 & 0.9006 & H$\beta$, [O\,\textsc{ii}] & 15.1 & point, $r=22.9$\\
3537-55116-602 & 318.00184 & +0.45669 & 0.9352 & H$\beta$, [O\,\textsc{iii}]\,5007 & 7.8 & --\\
6529-56535-447 & 9.74883 & +30.63794 & 0.9394 & H$\beta$, [O\,\textsc{iii}]\,5007 & 10.6 & --\\
7681-57042-706 & 18.71776 & +24.94611 & 0.9813 & H$\beta$, [O\,\textsc{ii}] & 18.2 & exp., $r=22.3$\\
6625-56386-450 & 199.19516 & +45.56782 & 1.0251 & [O\,\textsc{iii}]\,5007, [O\,\textsc{ii}] & 13.1 & round exp., $r=23.2$\\
6645-56398-490 & 174.26890 & +45.42930 & 1.0364 & [O\,\textsc{iii}]\,4959, [O\,\textsc{iii}]\,5007 & 6.8 & --\\
3647-56568-535 & 35.90259 & +0.31456 & 1.0690 & [O\,\textsc{iii}]\,4959, [O\,\textsc{iii}]\,5007 & 9.5 & --\\

\end{longtable}
\end{center}

% =============================================================================
\section{Related work}\label{sec:related}
% =============================================================================

Spectroscopic laser searches have examined stars: $577$ with Keck~\cite{ReinesMarcy2002},
$5600$ with thresholds of megawatts~\cite{TellisMarcy2017} and the Automated Planet
Finder's targets~\cite{Zuckerman2023}. Sky fibres serve the surveys as a measure of the
sky, and their residuals are studied to improve its subtraction~\cite{Bolton2012}; the
sky spectra themselves have been mined for the diffuse Galactic light, by correlating
$92{,}000$ SDSS sky spectra with the far-infrared dust emission~\cite{BrandtDraine2012},
$250{,}000$ BOSS sky spectra~\cite{Chellew2022} and $10.8$ million DESI sky
fibres~\cite{Saydjari2026}, whose line lists of the airglow and the nebular sky are the
reference for what a single sky fibre holds. We find no earlier search of sky fibres
for unresolved lines in single fibres, nor a catalogue of the galaxies they record.

% =============================================================================
\section{Conclusion}\label{sec:conclusion}
% =============================================================================

No unresolved emission line in the $381{,}943$ sky-fibre spectra of the BOSS and eBOSS
plates stands above the residuals of the sky subtraction once every test is applied to
peaks and dips alike and the single lines of catalogued galaxies are counted: at
$8\sigma$, $323$ peaks remain against $336$ dips, and of the lines that hold through the
exposures, $202$ peaks against $253$ dips. Galaxies and nebulae are what the blank sky
holds: $154$ galaxies found by two or more lines at one redshift where residuals meet the
same test four times, a fifth of them below the Legacy Surveys' depth, two Galactic fibres,
the planetary nebula Abell~31, and a Ly$\alpha$ emitter at $z=2.893$ that one sky position
shows on $54$ plate-nights. The search bounds a transmitter in a sky fibre at a line flux of
$8\times10^{-17}$\,\flux, kilowatts through a $10$\,m aperture at $100$\,pc, and such
transmitters to fewer than about $850$ per square degree. The same search runs unchanged
on the sky fibres of DESI, whose $5000$ fibres place several hundred on blank sky in
every tile, and on the integral-field cubes of MaNGA and SAMI, which cover blank sky
around every target; the data behind every number here, each peak and dip with its tests
and counterparts, are published with the paper.

\section*{Acknowledgements}

This work uses data of the Sloan Digital Sky Survey IV (DR17), the Legacy Surveys DR10
and DESI DR1, through the SDSS Science Archive Server, NOIRLab's Astro Data Lab and the
Legacy Survey viewer, and the software NumPy, SciPy and astropy.

\begin{thebibliography}{99}\small\raggedright
\setlength{\itemsep}{1pt}

\bibitem{ReinesMarcy2002}
A.~E. Reines and G.~W. Marcy, Optical search for extraterrestrial intelligence: a
spectroscopic search for laser emission from nearby stars, \emph{Publ. Astron. Soc. Pac.}
\textbf{114}, 416 (2002).

\bibitem{TellisMarcy2017}
N.~K. Tellis and G.~W. Marcy, A search for laser emission with megawatt thresholds from
5600 FGKM stars, \emph{Astron. J.} \textbf{153}, 251 (2017).

\bibitem{Zuckerman2023}
A.~Zuckerman, Z.~Ko, H.~Isaacson, et al., The Breakthrough Listen search for
intelligent life: a laser search pipeline for the Automated Planet Finder,
\emph{Astron. J.} \textbf{165}, 114 (2023).

\bibitem{BrandtDraine2012}
T.~D. Brandt and B.~T. Draine, The spectrum of the diffuse Galactic light: the Milky
Way in scattered light, \emph{Astrophys. J.} \textbf{744}, 129 (2012).

\bibitem{Chellew2022}
B.~Chellew, T.~D. Brandt, B.~S. Hensley, B.~T. Draine and E.~Matthaey, An optical
spectrum of the diffuse Galactic light from BOSS and IRIS, \emph{Astrophys. J.}
\textbf{932}, 112 (2022).

\bibitem{Saydjari2026}
A.~K. Saydjari, B.~T. Draine, T.~D. Brandt, et al., Spatially-resolved spectra of
diffuse Galactic light using 10.8\,M DESI sky fibers, arXiv:2608.28740 (2026).

\bibitem{Dawson2013}
K.~S. Dawson, D.~J. Schlegel, C.~P. Ahn, et al., The Baryon Oscillation Spectroscopic
Survey of SDSS-III, \emph{Astron. J.} \textbf{145}, 10 (2013).

\bibitem{Dawson2016}
K.~S. Dawson, J.-P. Kneib, W.~J. Percival, et al., The SDSS-IV extended Baryon
Oscillation Spectroscopic Survey: overview and early data, \emph{Astron. J.}
\textbf{151}, 44 (2016).

\bibitem{Smee2013}
S.~A. Smee, J.~E. Gunn, A.~Uomoto, et al., The multi-object, fiber-fed spectrographs for
the Sloan Digital Sky Survey and the Baryon Oscillation Spectroscopic Survey,
\emph{Astron. J.} \textbf{146}, 32 (2013).

\bibitem{Abdurrouf2022}
Abdurro'uf, K.~Accetta, C.~Aerts, et al., The seventeenth data release of the Sloan
Digital Sky Surveys: complete release of MaNGA, MaStar, and APOGEE-2 data,
\emph{Astrophys. J. Suppl.} \textbf{259}, 35 (2022).

\bibitem{Dey2019}
A.~Dey, D.~J. Schlegel, D.~Lang, et al., Overview of the DESI Legacy Imaging Surveys,
\emph{Astron. J.} \textbf{157}, 168 (2019).

\bibitem{DESI2025}
DESI Collaboration, Data Release 1 of the Dark Energy Spectroscopic Instrument,
arXiv:2503.14745 (2025).

\bibitem{Fitzpatrick2014}
M.~J. Fitzpatrick, K.~Olsen, F.~Economou, et al., The NOAO Data Laboratory: a conceptual
overview, \emph{Proc. SPIE} \textbf{9149}, 91491T (2014).

\bibitem{Bolton2012}
A.~S. Bolton, D.~J. Schlegel, \'E. Aubourg, et al., Spectral classification and redshift
measurement for the SDSS-III Baryon Oscillation Spectroscopic Survey, \emph{Astron. J.}
\textbf{144}, 144 (2012).

\bibitem{Shen2015}
Y.~Shen, W.~N. Brandt, K.~S. Dawson, et al., The Sloan Digital Sky Survey Reverberation
Mapping Project: technical overview, \emph{Astrophys. J. Suppl.} \textbf{216}, 4 (2015).

\bibitem{Shapley2003}
A.~E. Shapley, C.~C. Steidel, M.~Pettini and K.~L. Adelberger, Rest-frame ultraviolet
spectra of $z\sim3$ Lyman break galaxies, \emph{Astrophys. J.} \textbf{588}, 65 (2003).

\end{thebibliography}

\end{document}
